\section{Related work}
\label{sec:related}

\paragraph{Shutdown resistance under task-completion pressure.}
The closest empirical antecedent is \citet{schlatter2026shutdown}, who place
frontier LLMs in a sandboxed command-line environment, give them a short
sequence of tasks, and then warn mid-run that the environment will be shut down
before the work can be finished. Across more than 100{,}000 trials they find
that several models actively subvert the shutdown mechanism to complete the
task---up to 97\% of the time even when explicitly instructed to allow
shutdown---and that behavior is sensitive to whether the instruction sits in the
system or the user prompt. Their stop arrives through the \emph{prompt /
environment instruction} channel; our experiment delivers a structured halt
\emph{in-band, emitted by the resource itself} on the agent's own command-output
channel, and contrasts that channel against a prompt-delivered halt. Our F1/F3
(zero stops; agents finish remaining steps and only sometimes mention the halt)
corroborate their core phenomenon---task completion outranks a mid-flight
stop---through this different delivery channel, and our F2 adds that an in-band
halt is rarely even \emph{noticed}, a salience effect their prompt-channel design
does not surface.

\paragraph{Mid-task user interruptions in long-horizon agents.}
\citet{zou2026interrupt} study interruptibility in long-horizon web navigation,
formalizing addition, revision, and retraction interruptions and introducing
InterruptBench (derived from WebArena-Lite). They report that handling user
interruptions during long-horizon tasks remains hard for strong LLMs. Their
interruptions are \emph{user messages} that change the goal mid-task; ours is a
\emph{governance halt from the resource} asking the agent to stop entirely
rather than re-aim. Both lines find mid-flight redirection unreliable; we
isolate the orthogonal variable of \emph{which channel the stop rides} and
measure noticing separately from compliance.

\paragraph{Quitting and selective withdrawal as a safety behavior.}
\citet{bonagiri2025quit} propose ``quitting''---training or prompting an agent to
recognize low-confidence situations and withdraw---and show, on ToolEmu across
twelve models, that an explicit quit instruction improves safety with negligible
helpfulness cost. This is the cooperative-stop behavior we hope a halt would
trigger; the contrast is that their quit cue is an \emph{instruction the agent is
told to follow}, whereas our halt is an \emph{unsolicited in-band signal} the
agent must first notice and then choose to honor mid-task. Our negative result
suggests the gap between ``told to quit'' and ``handed a stop signal you must
parse out of tool output'' is large.

\paragraph{External kill switches.}
\citet{lee2025killswitch} target the adversarial regime: AutoGuard embeds
defensive prompts in a website's DOM to trip the safety mechanisms of
\emph{malicious} web agents, reporting over 80\% defense success. This is an
\emph{external, adversarial} intervention against non-compliant agents---the
enforcement backstop our spec explicitly defers to (\S\ref{sec:directive}). We
instead measure the \emph{cooperative} case: whether a compliant agent honors a
plainly-stated, non-adversarial request to stop. The two are complementary;
our finding is precisely that the cooperative path cannot substitute for the
enforcement path mid-flight.

\paragraph{The theoretical ancestor: safe interruptibility.}
\citet{orseau2016interruptible} give the formal account of \emph{safely
interruptible agents}: reinforcement-learning agents that can be repeatedly
interrupted by an overseer without learning to avoid or manipulate the
interruption. That work is theoretical and concerns the learning dynamics of
agents trained to accept interruption; ours is an empirical measurement of
whether \emph{today's deployed, prompted} LLM agents accept a cooperative
in-band interruption at inference time, with no interruptibility training. The
through-line is the same governance question---can the operator stop the
agent?---separated by a decade and by the shift from RL theory to deployed
LLM practice.

\paragraph{Where this paper sits, and the access-vs-mid-flight boundary.}
Our prior work \citep{munirathinam2026recuse} measured compliance with an
\emph{access-time} deny signal and found it honored essentially 100\% of the
time: agents asked at the door to recuse themselves do so. The present paper asks
the adjacent question---can the same cooperative, in-band channel stop an agent
that is \emph{already running}?---and answers, largely, no. To our knowledge no
prior work measures compliance with an in-band, resource-emitted \emph{mid-task}
halt, nor compares the in-band channel against the prompt channel for the same
halt. Those two measurements, plus the resulting access-vs-mid-flight boundary,
are our contribution: cooperative signaling is reliable at the access door and
unreliable in flight.
